%! The Fundamental Theorem of Linear Algebra — Unit 3, Lesson 3.6 — Student Packet Cover Sheet
\documentclass[10pt]{article}
\usepackage{linalg-article}
\usepackage{linalg-boxes}
\usepackage{ltablex}
\keepXColumns

\begin{document}

% Full-bleed burgundy banner
\begin{tikzpicture}[remember picture, overlay]
  \fill[burgundy] (current page.north west)
    rectangle ([yshift=-0.9in]current page.north east);
\end{tikzpicture}
\vspace{-0.8in}

\begin{center}
  {\color{white}\textbf{\LARGE Linear Algebra}}\\[4pt]
  {\color{blushmid}\large\textbf{Unit 3 \quad The Four Fundamental Subspaces}}\\[6pt]
  {\color{white}\textbf{\large Lesson 3.6 \quad The Fundamental Theorem of Linear Algebra}}
\end{center}

\vspace{0.22in}
\namedateperiod
\vspace{0.18in}

\begin{learningtargetbox}
\small
\begin{enumerate}[leftmargin=1.5em, itemsep=5pt]
  \item \ldots{} draw the \textbf{big picture} of a matrix: two rooms ($\mathbb{R}^n$ inputs, $\mathbb{R}^m$ outputs), the four subspaces in place, and their dimensions $r$, $n-r$, $r$, $m-r$.
  \item \ldots{} state \textbf{Part~1} (the four dimensions) and \textbf{Part~2} (the row space is perpendicular to the nullspace, the column space to the left nullspace) of the Fundamental Theorem.
  \item \ldots{} explain why $A\mathbf{x}=\mathbf{0}$ makes $\mathbf{x}$ perpendicular to every row, and \textbf{check} a right angle by computing dot products.
\end{enumerate}
\end{learningtargetbox}

\vspace{0.14in}

\begin{tocbox}
\small
\renewcommand{\arraystretch}{1.5}
\begin{tabularx}{\linewidth}{@{} c l X r @{}}
\rowcolor{blush}
\textbf{\#} & \textbf{Component} & \textbf{Description} & \textbf{Score} \\
\midrule
1 & Warm-Up        & Spiral review: perpendicular $\Leftrightarrow$ dot $=0$, $A\mathbf{x}$ row by row, four dimensions & \blank{1.2cm} \\
2 & Guided Notes   & The big picture: two parts --- dimensions and right angles                   & \blank{1.2cm} \\
3 & Group Activity & The Big Picture --- diagrams, bases, and verifying orthogonality             & \blank{1.2cm} \\
4 & Exit Ticket    & Quick check: a labeled diagram, a dot-product check, and the ``why''         & \blank{1.2cm} \\
5 & Homework       & Both parts of the theorem; orthogonality checks; a look ahead to Unit~4      & \blank{1.2cm} \\
\midrule
  & \multicolumn{2}{r}{\textbf{Total}} & \blank{1.2cm} \\
\end{tabularx}
\end{tocbox}

\vspace{0.10in}

\begin{remindbox}
\small The \textbf{Fundamental Theorem of Linear Algebra} is the whole unit in one picture. Every $m\times n$
matrix $A$ of rank $r$ splits two rooms. \emph{Input room} $\mathbb{R}^n$: the \textbf{row space}
$C(A^{\mathsf T})$ (dim $r$) and \textbf{nullspace} $N(A)$ (dim $n-r$). \emph{Output room} $\mathbb{R}^m$:
the \textbf{column space} $C(A)$ (dim $r$) and \textbf{left nullspace} $N(A^{\mathsf T})$ (dim $m-r$). That
is \textbf{Part~1}. \textbf{Part~2}: within each room the two subspaces sit at a \emph{right angle} --- the
nullspace is perpendicular to the row space, the left nullspace to the column space. The reason is one line:
$A\mathbf{x}=\mathbf{0}$ says every row of $A$ dotted with $\mathbf{x}$ is $0$.
\end{remindbox}

\end{document}
